% =====================================================================
% TEMPLATE — Scientific paper manuscript (methodology asset A-P11-tpl)
% =====================================================================
% Purpose: the OPTIONAL phase-P11 deliverable — a submittable scientific-journal
% manuscript that turns the finished case study (the P0–P10 artifacts) into a
% coherent paper. This is a JOURNAL-AGNOSTIC skeleton: a neutral `article`
% default that the paper-writer agent ADAPTS to a specific target journal.
%
% Fill every <PLACEHOLDER>; delete guidance comments. This file is meant to be
% BOTH standalone-compilable (it carries its own \documentclass ... \end{document}
% below) AND, if a project prefers, \input-able from a thin main.tex (in that
% case strip the preamble + \begin/\end{document} and keep only the body).
%
% Section skeleton (keep it stable; adapt ORDER/NAMING to the target journal):
%   Title + authors → Abstract → Keywords → Introduction (motivation +
%   explicit contribution bullets) → Related work (\input related_work.tex) →
%   Problem statement → Methods (Exact + Heuristic) → Experiments & results →
%   Discussion (honest caveats) → Conclusions & future work → References.
%
% Inputs this manuscript draws from (do NOT reproduce them — summarise + cite):
%   * P1 problem_description.tex   — problem statement, [RQxx] requirements
%   * P2 data_model.tex            — data / instance definitions
%   * P5 milp_formulation.tex      — the exact model (cite the detailed doc)
%   * P8 heuristic.tex             — the heuristic (cite the detailed doc)
%   * P9 ledger / P10 report.tex   — results, tables, gap convention, caveats
%   * P11 related_work.tex + references.bib — scholarly positioning (this phase)
%
% Discipline (inherited from the report + the P11 agents):
%   * No invented numbers — every result traces to the P9 ledger / P10 report.
%   * No invented citations — every \cite resolves in references.bib.
%   * Carry over the report's HONEST CAVEATS; do not overclaim.
% =====================================================================

% ---------------------------------------------------------------------
% DOCUMENT CLASS + CITATION STYLE — choose per TARGET JOURNAL.
%   * Neutral default below: `article` + natbib author–year. Good when no target
%     journal is given (state the assumption in the paper-writer's report).
%   * For a specific journal, SWAP this preamble for the publisher's class and
%     citation style, e.g.:
%       - Elsevier:  \documentclass{elsarticle}      (numeric, \usepackage-free bib)
%       - Springer:  \documentclass{sn-jnl}          (sn-mathphys / basic bst)
%       - INFORMS:   \documentclass{informs3}
%       - IEEE:      \documentclass{IEEEtran}         (IEEEtran bst, numeric)
%     Match the journal's CITATION FORMAT: numeric (\cite → [1]) vs author–year
%     (\citep/\citet). Keep the SAME \label/\cite keys regardless of style so the
%     bibliography (references.bib) is portable.
% ---------------------------------------------------------------------
\documentclass[11pt]{article}
\usepackage[utf8]{inputenc}
\usepackage{amsmath,amssymb}
\usepackage{graphicx}
\usepackage{booktabs}
\usepackage[round]{natbib}   % author–year default; switch to journal style
\usepackage{hyperref}

\title{<PAPER TITLE: concise, states the problem + the contribution>}

% Authors + affiliations: leave a clearly-marked TODO rather than inventing them.
\author{<AUTHOR NAMES>\\ <AFFILIATION(S)> \\ \texttt{<CONTACT EMAIL>}}
% % TODO authors/affiliations/funding — fill before submission
\date{<DATE or \today>}

\begin{document}
\maketitle

% ---------------------------------------------------------------------
\begin{abstract}
% 150–250 words (respect the journal's limit): the problem and why it matters,
% what is proposed (exact model + heuristic), the headline result (the single
% strongest evidence-backed finding), and the takeaway. No citations, no
% undefined acronyms. Numbers must match the results section exactly.
<ABSTRACT: problem + motivation, the methods proposed, the headline result,
the takeaway>
\end{abstract}

\noindent\textbf{Keywords:} <KEYWORD 1>; <KEYWORD 2>; <KEYWORD 3>; <...>
% ^ 4–6 keywords; prefer the journal's controlled vocabulary if it has one.

% =====================================================================
\section{Introduction}
% Funnel: (1) the real-world / scientific context and why the problem matters;
% (2) the specific problem and what makes it hard; (3) the gap in existing work
% (one sentence, expanded in Related work); (4) what THIS paper does.
% Keep it self-contained — a reader outside the sub-field should follow it.
<INTRODUCTION: motivation, the problem, the gap, and what this paper does>

% Explicit contribution bullets — the reviewer's checklist. Each is a concrete,
% defensible claim backed later by the paper (a model, an algorithm, an
% experimental finding). Do not list anything you cannot substantiate.
\paragraph{Contributions.}
\begin{itemize}
  \item <CONTRIBUTION 1: e.g. a formal model of the problem>.
  \item <CONTRIBUTION 2: e.g. an exact formulation and/or a heuristic>.
  \item <CONTRIBUTION 3: e.g. an experimental comparison establishing X>.
  \item <CONTRIBUTION 4 (optional): reusable artifacts / instances released>.
\end{itemize}

% One sentence on the paper's structure (optional; some journals omit it).
<OPTIONAL: "The remainder of the paper is organized as follows ...">

% =====================================================================
\section{Related work}
% Scholarly POSITIONING (phase P11 bibreviewer output): thematic groups of prior
% work + an explicit gap/positioning paragraph. Kept in a separate file so the
% bibreviewer owns it and the writer owns the rest.
\input{related_work.tex}

% =====================================================================
\section{Problem statement}
% State the problem precisely: the inputs (from the P2 data model), the decisions,
% the constraints/requirements (reference the [RQxx] ids from P1), and the
% objective. Formal but concise — enough to make the methods well-defined; the
% full data model lives in the cited data-model doc.
<PROBLEM STATEMENT: inputs, decisions, constraints/[RQxx] requirements, objective>

% =====================================================================
\section{Methods}

\subsection{Exact <EXACT METHOD NAME>}
% Summarise the exact model and CITE the detailed formulation doc rather than
% reproducing every constraint: the decision variables/objects, the objective,
% the governing constraint families, and the solver stack + run knobs. If space
% allows, show only the objective and 1–2 signature constraints.
<EXACT METHOD SUMMARY: variables/objects, objective, key constraint families,
solver stack + run knobs; cite the detailed formulation doc>

\subsection{Heuristic <HEURISTIC METHOD NAME>}
% Name the approach (and the literature family it belongs to, cited). Summarise
% its components as a short list; CITE the detailed algorithm doc for pseudocode.
\begin{enumerate}
  \item \textbf{<CONSTRUCTIVE STEP>} --- <requirement-respecting builder; note
        which hard requirements ([RQxx]) it enforces by construction>.
  \item \textbf{<OBJECTIVE EVALUATOR>} --- <the deterministic evaluator producing
        a full, checker-valid solution + score (the reusable core)>.
  \item \textbf{<SEARCH / METAHEURISTIC>} --- <the improvement loop: degree of
        freedom perturbed, accept rule, restarts>.
\end{enumerate}
% One sentence: both methods emit the SAME solution schema and are validated by
% the same checker, so the comparison is apples-to-apples.
<CLOSING: both methods emit the same solution schema, judged identically>

% =====================================================================
\section{Experiments and results}
% Setup FIRST (brief), then the results table + narrative.
% Setup: the benchmark (deterministic instances from the P2 generator, size grid),
% the harness (dual-logged, compliance-gated by the P3 checker), the ONE gap
% convention (defined once), and reproducibility stamps. Numbers as logged.
<EXPERIMENTAL SETUP: benchmark size grid, harness + compliance gate, the fixed
gap convention (state what a positive value means), reproducibility>

% Results table — one row per instance (or per size class). Keep numbers EXACTLY
% as in the P9 ledger / P10 report.
\begin{table}[t]
\centering
\caption{<TABLE CAPTION: what is compared, over which instances>}
\label{tab:results}
\begin{tabular}{l<COLUMN SPEC: one column per reported metric>}
\toprule
instance & <BASELINE> obj & <BASELINE> status & <BASELINE>\,s & <METHOD> obj & <METHOD>\,s \\
\midrule
<INSTANCE> & <obj> & <status> & <time> & <obj> & <time> \\
% ... one row per instance / size class
\bottomrule
\end{tabular}
\end{table}

% Narrative: the headline finding stated plainly — the gap across the benchmark,
% how many instances the exact method closed, behavior where it stalled (objective
% + speed factor), that every reported solution is checker-valid, runtime envelope.
<RESULTS NARRATIVE: headline gap; closed count; behavior on unclosed instances;
all solutions checker-valid; runtime envelope>

% =====================================================================
\section{Discussion}
% Interpret the results, then state the HONEST CAVEATS explicitly (carry them over
% from the P10 report — do not launder them out for publication). Typical:
%   (i)   benchmark size/coverage — correctness validation vs scaling study;
%   (ii)  was the exact method actually stressed (was a positive gap demonstrable)?
%   (iii) determinism/variance limits.
<DISCUSSION: interpretation + the honest caveats bounding what the results
do and do not establish>

% =====================================================================
\section{Conclusions and future work}
% 2–4 sentences: the evidence-grounded takeaway (which method, and why), then the
% concrete next steps (scaling study, parameter sweeps, richer instances, etc.).
<CONCLUSIONS: the evidence-grounded takeaway + concrete future work>

% =====================================================================
% Acknowledgements / funding / data-availability — optional, journal-dependent.
% \paragraph{Acknowledgements.} <... or delete>
% \paragraph{Data availability.} <instances/code release note, or delete>

% =====================================================================
% References — style must match the target journal (see the class comment above).
% All \cite keys resolve in references.bib (produced by the P11 bibreviewer).
\bibliographystyle{plainnat}  % swap for the journal's .bst (e.g. elsarticle-num)
\bibliography{references}

\end{document}
